% ZTC: original zeta, original shifted arithmetic family, retained theta data.
\section{Original zeta and the retained theta comparison}
\label{sec:original-zeta-theta-correction}

\paragraph{Correction and provenance.}
The statement that the Rodgers--Tao heat kernel was the programme's
original construction over the entire absolute base was not established.
THS1--31 begin with that classical analytic kernel and construct a
support-preserving coefficient lift. HTB1--35 establish its exact
comparison with the shifted Hurwitz family. Neither construction is a
derivation of the scalar kernel from the entire arithmetic geometry.
The distinction is measurable: THS22 already proves that all supported
zero times act identically on the actual theta orbit. The ambient
coefficient injection does not make that orbit injective.

The original arithmetic functions in this calculation remain
\begin{equation}\tag{ZTC1}
 \zeta(s)=\sum_{n=1}^{\infty}n^{-s},\qquad
 F_t(s)=\sum_{n=0}^{\infty}(n+1+t)^{-s},
 \qquad \operatorname{Re}s>1,\quad t>-1.
\end{equation}
Their continuations, endpoints, and exact shift generator are proved in
HTB2--5, with the programme source identified there by version, hash,
and original theorem labels. In particular $F_0=\zeta$,
$F_1=\zeta-1$, and $\partial_tF_t(s)=-sF_t(s+1)$.
The latter family comes from the supplied April 2026 shifted-Hurwitz
programme draft, not from a proposed identification of two time variables.
Brad Rodgers and Terence Tao's
\href{https://arxiv.org/abs/1801.05914v5}{original author source},
equations \texttt{phidef}, \texttt{htdef}, \texttt{hoz}, and
\texttt{sas}, supplies the auxiliary heat kernel used in THS and HTB.
These are cited inputs, not new discoveries. The following comparisons
are derived explicitly. Their complete HTB, THS, and CFD proofs accompany
this chapter.

\subsection{A direct theta sum for the original arithmetic family}

For real $t>-1$ and $x>0$, retain the actual shifted frequencies and put
\begin{equation}\tag{ZTC2}
 K_t(x)=\sum_{n=0}^{\infty}
               e^{-\pi(n+1+t)^2x}.
\end{equation}
This is a scalar analytic representation of the stated arithmetic
family. No assertion that it determines the full support lattice is
part of this definition. For fixed $a=1+t>0$, monotonicity of the
summand as a function of $n$ gives
\begin{equation}\tag{ZTC3}
 K_t(x)\le e^{-\pi a^2x}
              +\int_0^\infty e^{-\pi(a+y)^2x}\,dy
 \le 1+\frac1{2\sqrt x}.
\end{equation}
For $x\ge1$ the same bound after retaining the first exponential
gives $K_t(x)\le C_a e^{-\pi a^2x}$, since
$(a+n)^2-a^2\ge n^2$ and
$C_a=\sum_{n\ge0}e^{-\pi((a+n)^2-a^2)}<\infty$.
Consequently the following integral converges absolutely for
$\operatorname{Re}s>1$:
\begin{equation}\tag{ZTC4}
 \int_0^\infty x^{s/2-1}K_t(x)\,dx
   =\pi^{-s/2}\Gamma(s/2)F_t(s),\qquad
 F_t(s)=\frac{\pi^{s/2}}{\Gamma(s/2)}
                \int_0^\infty x^{s/2-1}K_t(x)\,dx.
\end{equation}
Indeed, with $\sigma=\operatorname{Re}s>1$, the sum of the
absolute termwise integrals equals
$\pi^{-\sigma/2}\Gamma(\sigma/2)\sum_{n\ge0}(n+a)^{-\sigma}$.
Fubini therefore applies. Substitution $y=\pi(n+a)^2x$ in each
term proves the displayed identity with its full Gamma and pi factors.
Continuation outside that domain means the meromorphic continuation
proved in HTB4, not an assertion that the unsubtracted integral converges
there.

There is also an exact inverse recovering the kernel, not only the
arithmetic function. For every real $c>1$,
\begin{equation}\tag{ZTC4a}
 K_t(x)=\frac1{4\pi i}\int_{c-i\infty}^{c+i\infty}
       \pi^{-s/2}\Gamma(s/2)F_t(s)x^{-s/2}\,ds.
\end{equation}
Here the integral is absolutely convergent. To prove the inversion
with its constant, set $\sigma=c/2>0$ and
$g(v)=\exp(\sigma v-e^v)$. This function and all its derivatives
are integrable on $\mathbb R$: at $-\infty$ they are bounded
by a constant times $e^{\sigma v}$, and at $+\infty$ by a
polynomial in $e^v$ times $g(v)$. The substitution $u=e^v$
identifies its Fourier transform with $\Gamma(\sigma-i y)$.
Repeated integration by parts therefore bounds that Gamma factor
by $C_N(1+|y|)^{-N}$ for every fixed $N$.
Fourier inversion gives
$e^{-u}=(2\pi i)^{-1}\int_{\sigma-i\infty}^{\sigma+i\infty}
\Gamma(w)u^{-w}\,dw$. One may verify this inversion directly
by first multiplying the Fourier integrand by $e^{-\epsilon y^2}$:
the resulting expression is the convolution of $g$ with the
Gaussian approximate identity, which converges to $g$; the stated
integrable transform bound permits dominated convergence as
$\epsilon\downarrow0$. Apply the identity to
$u=\pi(n+1+t)^2x$, then put $s=2w$. The resulting factor is
$1/(4\pi i)$. Finally sum in $n$. Interchange is justified by
$\sum(n+1+t)^{-c}<\infty$ and the integrable Gamma bound.
This proves ZTC4a. ZTC4 and ZTC4a are mutual recovery formulas
for the specified original arithmetic family and its scalar
theta sums; they do not recover a support lattice from a scalar
function.

Differentiation in $t$ is locally uniform for $t$ in compact subsets
of $(-1,\infty)$ and $x$ in compact subsets of $(0,\infty)$, by
Gaussian majorants for a polynomial times $e^{-c n^2}$.
It gives the exact formula
\begin{equation}\tag{ZTC5}
 \partial_tK_t(x)=-2\pi x\sum_{n=0}^{\infty}
              (n+1+t)e^{-\pi(n+1+t)^2x}.
\end{equation}
For $\operatorname{Re}s>0$, integration of its absolute value is
bounded by a constant times $\sum(n+a)^{-\sigma-1}$, locally
uniformly in $t$. More precisely its Mellin transform is
\begin{equation}\tag{ZTC6}
 \int_0^\infty x^{s/2-1}\partial_tK_t(x)\,dx
 =-2\pi^{-s/2}\Gamma(s/2+1)F_t(s+1)
 =-s\pi^{-s/2}\Gamma(s/2)F_t(s+1).
\end{equation}
Thus the original arithmetic shift generator is represented directly,
without substituting the heat generator for it. The stated equality
extends meromorphically by HTB4--5.

The separated sheet and its entire missing summand are visible before
any transform:
\begin{equation}\tag{ZTC7}
 K_1(x)=K_0(x)-e^{-\pi x},\qquad
 D_t(s)=\zeta(s)-F_t(s),\qquad D_1(s)=1,
\end{equation}
\begin{equation}\tag{ZTC8}
 \pi^{-s/2}\Gamma(s/2)D_t(s)
   =\int_0^\infty x^{s/2-1}\bigl(K_0(x)-K_t(x)\bigr)\,dx
       \quad(\operatorname{Re}s>1).
\end{equation}
ZTC7 follows by deleting the $n=1$ summand of $K_0$; ZTC8 follows
by subtracting two absolutely convergent integrals. At $t=1$ the
integral is exactly $\pi^{-s/2}\Gamma(s/2)$, so its inverse returns
the actual constant $1$, including its amplitude and support when lifted.

\subsection{Original zeta with both Mellin endpoints retained}

At $t=0$, HTB6 proves
\begin{equation}\tag{ZTC9}
 K_0(x)=x^{-1/2}K_0(1/x)+\tfrac12x^{-1/2}-\tfrac12.
\end{equation}
Split ZTC4 at $x=1$, apply this identity only to the integral over
$(0,1)$, and substitute $y=1/x$ there. Initially
$\operatorname{Re}s>1$, and direct integration gives
\begin{equation}\tag{ZTC10}
 \zeta(s)=\frac{\pi^{s/2}}{\Gamma(s/2)}
 \left\{
 \frac1{s-1}-\frac1s
 +\int_1^\infty K_0(x)
             \left(x^{s/2-1}+x^{-(s+1)/2}\right)\,dx
 \right\}.
\end{equation}
The two separate rational terms are the transforms of the two separate
terms $x^{-1/2}/2$ and $-1/2$. The integral on $[1,\infty)$ is
entire in $s$, since exponential decay bounds every logarithmic
derivative on compact subsets of the $s$-plane. This proves the
meromorphic continuation of the identity and keeps both endpoint
contributions. Neither has been declared zero or absorbed into an
unrecorded constant.

For the original cosine heat kernel, retain
\begin{equation}\tag{ZTC11}
\begin{gathered}
 A(s)=\frac{s(s-1)}2\pi^{-s/2}\Gamma(s/2),\qquad
 s(z)=\frac12+\frac{iz}{2},\qquad z(s)=-2i(s-1/2),\\
 H_0(z)=\frac{s(z)(s(z)-1)}{16}
           \pi^{-s(z)/2}\Gamma(s(z)/2)\zeta(s(z)),\\
 \zeta(s)=\frac{16\pi^{s/2}}
                  {s(s-1)\Gamma(s/2)}H_0(z(s)).
\end{gathered}
\end{equation}
The proof is HTB6--10: the exact Mellin substitution and two
integrations by parts yield the displayed multiplier, with all
endpoint limits justified in the original convergence half-plane.
The last equality is an identity of meromorphic functions. It is
not a licence to replace the original zeta calculation by an entire
function calculation while dropping its multiplier.

\subsection{The boundary is part of the receiving object}

Let $\mathcal M_s$ and $\mathcal M_z$ be the complex vector spaces
of meromorphic functions on the respective complex planes. Retain
the linear isomorphism, already proved in HTB11,
\begin{equation}\tag{ZTC12}
 Uf(z)=\frac{A(s(z))}{8}f(s(z)),\qquad
 U^{-1}G(s)=\frac{8G(z(s))}{A(s)}.
\end{equation}
Define the joint map on the entire two-component space by
\begin{equation}\tag{ZTC13}
 \mathcal V:\mathcal M_s\oplus\mathcal M_s
          \longrightarrow\mathcal M_z\oplus\mathcal M_z,
 \qquad (F,D)\longmapsto\bigl(U(F+D),UD\bigr).
\end{equation}
It has the exact inverse
\begin{equation}\tag{ZTC14}
 \mathcal V^{-1}(H,J)
    =\left(\frac{8(H-J)\circ z}{A},\frac{8J\circ z}{A}\right).
\end{equation}
Substitution proves both compositions equal the identity; all entries
are meromorphic functions, so no undefined product of separate
singular point values is involved. On the actual arithmetic family,
\begin{equation}\tag{ZTC15}
\begin{gathered}
 J_t=UD_t,\qquad
 \mathcal V(F_t,D_t)=(H_0,J_t),\\
 F_t(s)=\frac{16\pi^{s/2}}
              {s(s-1)\Gamma(s/2)}
              \bigl(H_0(z(s))-J_t(z(s))\bigr),\\
 J_1(z)=\frac{A(s(z))}{8},\qquad
 J_t(i)=-\frac t8.
\end{gathered}
\end{equation}
For the last equality use HTB5:
$D_t(0)=\zeta(0)-F_t(0)=t$, $A(0)=-1$, and $z(0)=i$.
In particular the retained pair determines the real parameter $t$
by $t=-8J_t(i)$. All $t$ give the same first component $H_0$.
Discarding $J_t$ therefore identifies different actual arithmetic
inputs. This is an explicit failure of injectivity on the actual
programme family, not merely an example in an unrelated vector space.
On the unrestricted space the summation map $(G,J)\mapsto G+J$
has exactly the kernel $\{(B,-B):B\in\mathcal M_z\}$.

For the complete support carrier, use the common label on the
direct sum:
\begin{equation}\tag{ZTC16}
 G_L(\mathcal M_s\oplus\mathcal M_s)
   \xrightarrow{\mathcal V^\uparrow}
 G_L(\mathcal M_z\oplus\mathcal M_z),\qquad
 ((F,D),\lambda)\longmapsto(\mathcal V(F,D),\lambda).
\end{equation}
Here $G_L(V)=\{(0,\lambda):\lambda\in L\}\cup(V\times\{1_L\})$
with its original joins and meets. Additivity, compatibility with
the $G_L(\mathbb C)$ action, carrier membership, and the inverse
follow componentwise from linearity and ZTC14, as explicitly proved
in CFD11--12. This is a bijection for every bounded distributive $L$.
It does not silently replace the common label by independent labels
on two product factors. In particular $e=(0,1_L)$ and
$\tau=(0,0_L)$ remain distinct when $L$ is nontrivial.

\subsection{Local zeros, poles, and the original trace}

The earlier conversational assertion that the singular module map
had not been established was also too broad: HTB13 already proves
$A^{-1}\mathcal O\xrightarrow{U}\mathcal O$ with the affine
coordinate change. CFD1--10 now make the comparison with the
\emph{original} regular lattice and every finite jet explicit.
At a finite point let $h=s-s_0$ and $A=h^k a(h)$, $a(0)\ne0$.
Retain $I=A^{-1}\mathcal O=h^{-k}\mathcal O$ as an embedded
submodule of the meromorphic germs, and $J=I\cap\mathcal O$.
For nonzero $f\in I$ the original order is
\begin{equation}\tag{ZTC17}
 \operatorname{ord}_{s_0}f
 =\dim_{\mathbb C}(I/\mathcal O f)
  -\dim_{\mathbb C}(I/J)
  +\dim_{\mathbb C}(\mathcal O/J).
\end{equation}
This is proved by the power-of-$h$ bases in CFD6--7, without
identifying the reference lattices. At a trivial zero $s_0=-2m$,
$I=h\mathcal O$, $I/\mathcal O\zeta=0$ and
$\mathcal O/I$ has dimension $1$. At $s_0=1$,
$I=h^{-1}\mathcal O$, $I/\mathcal O\zeta=0$ and
$I/\mathcal O$ has dimension $1$. Their signs in ZTC17 give,
respectively, $+1$ and $-1$. At $s_0=0$, $A(0)=-1$ is a unit.
CFD4 provides the invertible triangular matrices recovering all
Laurent coefficients with the exact factors $8$ and $i/2$.
Keeping only the zero algebra of the completed section would give
zero in both cancelled cases; keeping these comparison modules
restores the actual divisor.

The earlier conversational explanation of the fibres requires a
further correction. The map
\begin{equation}\tag{ZTC17a}
 \mathcal O/h\mathcal O\longrightarrow
 h\mathcal O/h^2\mathcal O,\qquad [u]\longmapsto[hu]
\end{equation}
is an isomorphism of $\mathcal O$-modules. It is well-defined because
$h(h\mathcal O)=h^2\mathcal O$, surjective by the definition of
$h\mathcal O$, and injective because $hu\in h^2\mathcal O$ implies
$u\in h\mathcal O$ in the local integral domain. Its inverse is
$[hu]\mapsto[u]$. Consequently the two abstract fibres alone do
not account for the loss. What must be retained is their specified
embedding and the original reference lattice: multiplication by
$A=h^{-1}a(h)$ identifies $I=h\mathcal O$ with $\mathcal O$, but
the length-one quotient $\mathcal O/I$ records the cancelled
original zero. Replacing the entire marked comparison by the
completed section's zero quotient discards precisely that datum.
Likewise multiplication by nonzero meromorphic $A$ is an
automorphism of the meromorphic sheaf. It is not an automorphism
of the fixed holomorphic lattice at its zeros and poles. These
category statements do not contradict one another.

Logarithmic differentiation gives the original trace identity
\begin{equation}\tag{ZTC18}
 \frac{\zeta'(s)}{\zeta(s)}
 =-2i\frac{H_0'(z(s))}{H_0(z(s))}
 -\frac1s-\frac1{s-1}+\frac12\log\pi-\frac12\psi(s/2).
\end{equation}
It first holds away from zeros and poles and then meromorphically;
$\psi=\Gamma'/\Gamma$. For a positively oriented finite contour
$\gamma$ avoiding all zeros and poles of the functions and all
singularities of every displayed integrand, and a holomorphic test $g$
on and inside it, multiplication by $g$ and integration gives
\begin{equation}\tag{ZTC19}
\begin{split}
 \frac1{2\pi i}\int_\gamma g\,\frac{\zeta'}\zeta\,ds
 ={}&\frac1{2\pi i}\int_\gamma
           g(s)(-2i)\frac{H_0'(z(s))}{H_0(z(s))}\,ds\\
 &-\frac1{2\pi i}\int_\gamma g(s)
   \left(\frac1s+\frac1{s-1}-\frac12\log\pi
                         +\frac12\psi(s/2)\right)ds.
\end{split}
\end{equation}
At every zero or pole, write a meromorphic function as
$(s-s_0)^r$ times a holomorphic unit. Its logarithmic derivative
then has residue $r$. The residue theorem proves that ZTC19
counts the original zeros with multiplicity and the original poles
with negative multiplicity, including every trivial zero inside
the contour. No infinite-contour limit or decay assumption is
hidden here. Extending this finite identity to a particular infinite
trace requires the actual convergence estimates of that trace.

\subsection{Two different trajectories through the same original zeta}

The original-$s$ function associated with the classical heat path is
\begin{equation}\tag{ZTC20}
 f_h^{\rm heat}(s)=\frac{16\pi^{s/2}}
                         {s(s-1)\Gamma(s/2)}H_h(z(s)),
 \qquad f_0^{\rm heat}=\zeta.
\end{equation}
CFD18--21 give its full receiving operator and prove
\begin{equation}\tag{ZTC21}
 \left.\partial_hf_h^{\rm heat}(-2)\right|_{h=0}=0,
 \qquad
 \left.\partial_tF_t(-2)\right|_{t=0}=-\frac16.
\end{equation}
For completeness, the heat equation and chain rule give
$\partial_hf=(Af)''/(4A)$. At $h=0$, $A\zeta$ is holomorphic
at $-2$ whereas $1/A$ has a simple zero there, giving the first
value. HTB5 gives the second as $2\zeta(-1)=-1/6$.
The discrepancy holds on the actual initial arithmetic function.
CFD21 further proves that all the zeros at $-2m$ remain simple
along the real heat path, while its residue at $1$ varies strictly;
the actual Hurwitz path has constant residue $1$ and values
$-B_{2m+1}(1+t)/(2m+1)$ at those points.

Thus the valid original-$\zeta$ receiver consists of the original
functions, their exact transform and inverse, the retained boundary
component, the embedded local comparison modules, and all support
labels. The scalar theta orbit by itself is a smaller observable.
These proofs neither exhibit an off-line zero of the original
$\zeta$ nor rule one out. They establish precisely which input and
which deformation a subsequent estimate is estimating.
